% Generated draft: requires mathematical and source review.
\documentclass[11pt]{article}
\usepackage[margin=25mm]{geometry}
\usepackage{fontspec}
\setmainfont{Latin Modern Roman}
\usepackage{amsmath,amssymb,mathtools,unicode-math}
\setmathfont{Latin Modern Math}
\usepackage{enumitem,needspace,xurl,hyperref}
\setlength{\parindent}{0pt}
\setlength{\parskip}{5pt}
\setlength{\emergencystretch}{3em}
\newcommand{\N}{\mathbb{N}}
\newcommand{\Z}{\mathbb{Z}}
\newcommand{\Q}{\mathbb{Q}}
\newcommand{\R}{\mathbb{R}}
\newcommand{\C}{\mathbb{C}}
\newcommand{\F}{\mathbb{F}}
\newcommand{\A}{\mathbb{A}}
\newcommand{\PP}{\mathbb{P}}
\newcommand{\E}{\mathbb{E}}
\newcommand{\eps}{\varepsilon}
\newcommand{\dd}{\,\mathrm{d}}
\begin{document}
\section*{UnsolvedMath ID: 20001413}
\textbf{Source problem code:} catalogue entry
\subsection{Definitions and mathematical statement}
Let $T$ denote a specified dispersing-billiard dynamical system, with its invariant probability measures $\mathcal M_T$. The source does not specify whether $T$ is the collision map or the billiard flow, the table geometry, dimension, horizon, or treatment of grazing singularities. For the map interpretation, define the variational entropy
\[
 h_*=\sup_{\nu\in\mathcal M_T}h_\nu(T).
\]
A measure of maximal entropy (MME) is $\mu_*\in\mathcal M_T$ with $h_{\mu_*}(T)=h_*$. The source asks whether an MME exists for dispersing billiards, what ``physical properties'' it has, and how it relates to periodic orbits. It does not prescribe uniqueness, a reference volume for physical basins, a mixing property, or a particular equidistribution claim.

For a precise possible periodic-orbit target, let $P_n$ be a finite set of regular period-$n$ points (with the convention exact period versus $T^nx=x$ to be specified), invariant under $T$, and let
\[
 \nu_n=\frac1{|P_n|}\sum_{x\in P_n}\delta_x.
\]
One may ask for growth of $|P_n|$ and weak convergence $\nu_n\to\mu_*$. Counting points versus primitive orbits, weighting each orbit equally versus by its period, and map-time versus flow-time sampling are separate choices left open in the original.

The supplied machine summary reports that, for standard planar finite-horizon Sinai billiard maps, current work gives a unique MME for every table, and that Carrand obtains periodic-orbit equidistribution under ``negligible singularities'' and ``sparse recurrence''. The record supplies neither the precise class and regularity of tables nor definitions of those two extra hypotheses or a full citation. These are attributed literature comments with their stated restrictions, not independent present-day assessments or assertions for every dispersing billiard.

The aid also announces an alternative abstract sufficient criterion for discontinuous maps: regular periodic sets must have full exponential growth at $h_*$; the number of points crowding any finite-partition cylinder must be subexponential; and periodic measures must be uniformly tight away from both the discontinuity set and boundaries of every fixed finite join of the partition. For a finite partition $\mathcal P$, the relevant joins are $\bigvee_{j=0}^{m-1}T^{-j}\mathcal P$. The reported conclusion is convergence of uniform periodic-point measures to the unique MME, without assuming that $\mathcal P$ is generating or that metric entropy is upper semicontinuous. Here uniqueness is part of the target framework; it is not manufactured by the definition of an MME. The precise ambient space, periodic-set selection, cylinder bound, neighborhood-tightness quantifiers, and other technical hypotheses are not quoted in the record. Accordingly this remains an attributed criterion description, rather than a fully specified new theorem or a verified billiard application.

In particular an MME need not be declared a physical measure merely by attaining maximal entropy. If ``physical'' is meant in the technical sense of a positive-reference-volume basin of empirical convergence, that reference measure and the relevant observable class must be specified separately.
\subsection{Short English statement}
The question asks about maximal-entropy measures for dispersing billiards and their connection to periodic orbits. The supplied summary reports existence and uniqueness in a specific Sinai-billiard setting and conditional equidistribution results. It also describes an abstract periodic-orbit criterion, but does not provide its full hypotheses. The map/flow setting and the meaning of physical properties remain open formulation choices.
\subsection{Sources}
\begin{itemize}
\item Local UnsolvedMath registry, permanent ID 20001413, exact source code catalogue entry; canonical \texttt{aim-dynamical-systems-notes.json}, \texttt{notes[70]}, source item 28, tag ``problem'', workshop \emph{Stochastic methods for non-equilibrium dynamical systems}; \url{https://aimath.org/pastworkshops/noneqdynsysproblems.pdf}.
\item The original spells ``billards''; spelling is normalized without changing the question. No proposer or separate bibliography is supplied. The summary names Carrand for conditional periodic-orbit equidistribution but gives no title, year, theorem number or full reference, and does not identify the authors of the claimed every-table unique-MME result. Original remarks and literature field are empty.
\item Supplied metadata: PARTIAL-PROGRESS, registry status ``partially\_solved'', statement status ``exact''; attempt 1, result ``partial\_result'', problem status at run ``partially\_solved'', recorded rigor ``proved'', novelty confidence ``low'', candidate type ``theorem''. The machine aid is explicitly not peer reviewed and requires independent verification before citation. No external source was visited.
\end{itemize}
\end{document}
